\chapter*{本册导读：生命怎样保存和改写信息}
\addcontentsline{toc}{chapter}{本册导读}
\markboth{本册导读}{}

前两册回答的是“物质怎样组成生命”：元素搭成分子，分子组成细胞，细胞里跑着一张代谢网络。但生命还有另一件更不可思议的事——它会\textbf{记住}。一颗受精卵“记得”怎样长成完整的身体，你“记得”父母的样子，一个物种“记得”亿万年环境变迁留下的印记。

这些记忆都写在同一种分子里：DNA。第三册讲的就是生命保存、复制、调节和改写信息的故事。第八篇从豌豆花出发，看孟德尔怎样把遗传变成可计算的规律，再看科学家怎样一步步证明 DNA 就是遗传物质。第九篇钻进分子机制：几十亿个字母怎样被抄写、转录、翻译成蛋白质，基因又怎样被打开、关小、调大——以及为什么“一个基因决定一个性状”几乎总是错的。第十篇把时间拉长到亿万年：突变、选择、漂变怎样写成进化史。第十一篇回到当下：测序、基因编辑、基因检测、转基因正在怎样改变医学、农业和我们每个人，以及——哪些问题至今没有答案。

这一册最重要的态度，请提前带在身上：基因很重要，但基因不是命运。性状来自基因、环境、发育历史和随机过程的共同作用。读懂这一点，比记住任何名词都更能保护你——保护你不被“基因算命”收割，也不被“基因万能”或“基因恐怖”的极端故事吓倒。
